\documentclass[10pt]{article}
\usepackage[utf8]{inputenc}
\usepackage[T1]{fontenc}
\usepackage{amsmath}
\usepackage{amsfonts}
\usepackage{amssymb}
\usepackage[version=4]{mhchem}
\usepackage{stmaryrd}
\usepackage{bbold}
\usepackage{graphicx}
\usepackage[export]{adjustbox}
\graphicspath{ {./images/} }

\begin{document}
Тогда неприводимое представление группы $G$, соответствующее орбите $\Omega \subset \mathfrak{g}^{*}$, при ограничении на $H$ разлагается на неприводимые компоненты, соответствующие тем орбитам $\omega \in \mathfrak{h}^{*}$, к-рые лежат в $p(\Omega)$, а представление группы $G$, индуцированное неприводимым представленуем группы $H$, соответствующим орбите $\omega \subset \mathfrak{h}^{*}$, разлагается на неприводимые компоненты, соответствующие тем орбитам $\Omega \subset g^{*}$, к-рые имеют непустое пересечение с прообразом $p^{-1}(\omega)$. Из этих результатов вытекают два следствия: если неприводимые представления $T_{i}$ соответствуют орбитам $\Omega_{i}, i=1,2$, то тензорное пропзведение $T_{1} \otimes T_{2}$ разлагается на неприводимые компоненты, соответствующие тем орбитам $\Omega$, к-рые лежат в арифметич. сумме $\Omega_{1}+\Omega_{2}$; квазирегулярное представление группы $G$ в пространстве функций на $G / H$ разлагается на неприводимые компоненты, соответствующие тем орбитам $\Omega \subset \mathfrak{g}^{*}$, для к-рых образ $p(\Omega) \subset \mathfrak{h}^{*}$ содержит нуль.

III. Т е о р и я х а p a к т е о в. Для характеров неприводимых представлений (как обобщенных функциї на групше) предложена [2] следующая универсальная формула:

\[
\chi(\exp X)=\frac{1}{p(X)} \int_{\Omega} e^{2 \pi i\langle F, X\rangle} \boldsymbol{\beta}(F),
\]

где $\exp : \mathfrak{g} \rightarrow G-$ экспоненциальное отображение алгебры Ји $\mathfrak{g}$ в группу $G, p(X)$ - квадратный корень из плотности инвариантной меры Хаара на $G$ в канонич. координатах, $\beta$ - форма объема на орбите $\Omega$, связанная с канонической 2 -формой $B_{\Omega}$ соотношением $\beta=$ $=\frac{B_{\Omega}^{k}}{k !}, \quad k=\frac{1}{2} \operatorname{dim} \Omega$. Эта формула справедлива для нильпотентных групп, разрешимых групп типа 1 , компактных групп, дискретной серии представлениї полупростых действительных групп и основной сериі представлений комплексных полупростых групп. Для нек-рых вырожденных серий представлений $\mathrm{SL}(3, \mathbb{R})$ формула неверна. Из формулы $(*)$ получается простая формула для вычисления инфинитезимального характера неприводимого представления $T_{\Omega}$, соответствующего орбите $\Omega$, а именно: каждому оператору Лапласа $\Delta$ на $G$ может быть сопоставлен $A d^{*} G$-инвариантный многочлен $P_{\Delta}$ на $\mathfrak{g}^{*}$ так, что значение инфинитезимального характера представления $T_{\Omega}$ на элементе $\Delta$ в точности равно значению $P_{\Delta}$ на $\Omega$.

IV. K о н с т р у к ц и ю неприводимого унитарного представления группы $G$ по ее орбите $\Omega$ в коприсоединенном представлении можно рассматривать как операцию квантования нек-рой гамильтоновой системы, для к-рой $\Omega$ играет роль фазового пространства, а $G$ роль многомерного некоммутативного времени (или группы симметрий). При әтом $G$-орбиты в коприсоединенном представлении - это все $G$-однородные симплектич. многообразия, допускающие квантование. Таким образом, вторую основную гипотезу можно переформулировать так: каждая әлементарная квантовая система с временем (или группой симметрий) $G$ получается квантованием из соответствующей классич. системы (см. [2]).

Обнаружена также связь О. м. с теорией вполне интегрируемых гамильтоновых систем (см. [11]).

Лum.: [1] К и р и л Л о в А. А., "Успехи матем. наук», 1962, т. 17 , в. 4 , с. $57-110 ;$ [2] е г о же, Элементы теории представлений, 2 изд., М., 1978; [3] Д и к с м ь е жћ., Универсальные обертыва юmие алгебры, пер. с франц., M., 1978; [4] $\mathbf{S}$ i m m s D. J., W o o d h o u s e N. M. J. Lectures on geometric guantization,
B. - Hdlb. - N. Y., 1976; [5] A u s l a n d e r L., K o s t a n t B., "Invent. math.», 1971. v. 14, p. 255-354; [6] M o o r e C.C., "Ann. Math.", 1965 , v. 82 , № 1, p. $146-82$; [7] R o t h s c h i ld L. P. , W o l f J. A., "Ann. Sci. Ecole norm. supér. Sér.4",1974, v. 7 , N 2, p. $155-74$; [8] Représentations des groupes de Lie ré-

\includegraphics[max width=\textwidth, center]{2023_07_08_d251f1b52d3006866f1cg-1}
1979, т: 249, № 3, c. 525-28; [10] K i r i i 1 o v A. A., Infinite dimensional groups, their representations, orbits, invariants, в сб.: Proc. Int. Congr. Math., Helsinki, 1978; [11] R e ym a n A. G., Se m e n o v - T i a n - S h a n s k y M. A., «In-
vent. math.", 1979 , v. $54, \mathcal{N}_{1} 1$ p. $81-100$. A. A. Кирuлsos. ОРБИТА точки $x$ относительно группы $G$, действующей на множестве $X$ (слева),- множество

\[
G(x)=\{g(x) \mid g \in G\}
\]

Множество

\[
G_{x}=\{g \in G \mid g(x)=x\}
\]

является подгруппої в $G$ и наз. с т а б и ли з а т о-

\includegraphics[max width=\textwidth, center]{2023_07_08_d251f1b52d3006866f1cg-1(1)}
$\begin{array}{llll}\text { т о ч к и } x & \text { о т но с и т е л ь н о } & G \text {. Отображение }\end{array}$ $g \mapsto g(x), g \in G$, индуцирует биекцию между $G / G_{x}$ и орбитой $G(x)$. О. любых двух точек из $X$ либо не пересекаются, либо совпадают; иначе говоря, О. определяют разбиение множества $X$. Фактормножество по отношению эквивалентности, определенному этим разбиением, наз. п р о с т р а н с т вом о р б и т, или ф а к т о рм н о же с т в о м $X$ по $G$, и обозначается $X / G$. Conoставление каждой точке ее О. определяет канонич. отображение $\pi_{X, G}: X \rightarrow X / G$. Стабилизаторы точек из одной O. сопряжены в $G$, точнее $G_{g(x)}=g G_{x} g^{-1}$. Если в $X$ имеется только одна 0 ., то $X-$ о д н о р о д н о е п р о с т р а н с т в о г р у п п ы $G$; говорят также, что $G$ де й с т в у е т на $X$ т р а нз и т и в но. Если $G-$ топологич. группа, $X$ - топологич. пространство и действие непрерывно, то $X / G$ обычно снабжается топологией, в к-рой множество $U \subset X / G$ открыто в $X / G$ тогда и только тогда, когда множөство $\pi_{X, G}^{-1}(U)$ открыто в $X$.

П р и ме р ы. 1) Пусть $G$ - группа поворотов плоскости $X$ вокруг фиксированной точки $a$. Тогда О.әто всевозможные окружности с центром в $a$ (в том числе и сама точка $a$ ). 2) Пусть $G$ - группа всех невырожденных линейных преобразований конечномерного действительного векторного пространства $V, X$ - множество всех симметрич. билинейных форм на $V$, а действие $G$ на $X$ определено формулой $(g f)(u, v)=f\left(g^{-1}(u), g^{-1}(v)\right)$ для любых $u, v \in V$. Тогда O. группы $G$ на $X$-множество форм, имеющих фиксированный ранг и сигнатуру.

Пусть $G$ - вещественная группа Ли, гладко действуюгая на дифференцируемом многообразии $X$ (см. Jи группа преобразований). Для любой точки $x \in X$ орбита $G(x)$ является погруженным подмногообразием в $X$, диффеоморфным $G / G_{x}$ (диффеоморфизм индуцирован отображением $g \mapsto g(x), g \in G)$. Әто подмногообразие не обязательно замкнуто в $X$ (не обязательно вложено). Классич. примером служит «обмотка тора», то есть О. действия аддитивной группы $\mathbb{R}$ на торе

\[
T^{2}=\left\{\left(z_{1}, z_{2}\right)\left|z_{i} \in \mathbb{C},\right| z_{i} \mid=1, i=1,2\right\},
\]

заданного формулой

\[
t\left(z_{1}, z_{2}\right)=\left(e^{i t_{z_{1}}}, e^{i \alpha t_{z_{2}}}\right), t \in \mathbb{R},
\]

где $\alpha$ - иррациональное действительное число; замыкание такой $О$. совпадает с $T^{2}$. Если $G$ компактна, то все 0 . являются вложенными подмногообразиями.

Если $G$ - алгебраич. группа, $X$ - алгебраич. многообразие над алгебраически замкнутым полем $k$, a действие регулярно (см. Алгебрацческая группа преобразований), то любая орбита $G(x)$ является гладким алгебраич. многообразием, открытым в своем замыкании $\overline{G(x)}$ (в топологии Зариского), причем в $\overline{G(x)}$ всегда содержится замкнутая О. группы $G$ (см. [5]). В этом случае морфизм $G \rightarrow G(x), g \mapsto g(x)$, индуцирует изоморфизм алгебраич. многообразий $G G_{x}$ и $G(x)$ тогда и только тогда, когда он сепарабелен (это условие всегда выполнено, если $k$ - поле нулевой характеристики). О. максимальной размерности образуют открытое в $X$ множество.

Описание структуры О. для данного цействия обычно сводится к указанию в каждой О. нек-рого единствен-


\end{document}